This thesis examines whether a practical Linux CLI assistant can be built
with small local LLMs under offline and hardware constraints.
The thesis proposes shAI, a modular RAG-based prototype composed of a
command-vs.-query classifier, a document retriever,
and a summarizer that generates answers from retrieved passages.

Tested on the \texttt{classifier-100} dataset, the classifier achieved up to
88\% accuracy with a 2.9s latency. For retrieval on the \texttt{sample-100}
collection, a baseline document-level approach reached high top-5 hits
(86.3\%–-92.0\%) and MRR (0.724–-0.806), but suffered from massive context bloat
up to 120K tokens. To resolve this, passage-level retrieval and document-level
retrieval with passage-level reranking were implemented. Evaluated on datasets
up to \texttt{sample-1000}, these optimized approaches drastically compressed
output sizes to ~5K tokens, maintaining high precision with latencies under 2s
for search and under 10s for full RAG generation.

The results demonstrate that retrieval-first, passage-level architectures are
highly practical for local CLI assistance, whereas agentic open-ended
approaches remain unreliable on small models. The thesis provides an open-source
prototype and empirical benchmarks to guide future work on constrained task
design.
